\chapter{Encephalopathy}
\index{Encephalopathy}

\section*{Definition and Overview}
Encephalopathy is a general term for any disease or damage that affects how the brain works, causing altered mental state, confusion, and changes in behaviour or consciousness. It is not a single disease but a syndrome with many causes, including liver failure, kidney failure, infections, lack of oxygen, toxins, and metabolic disturbances. Encephalopathy ranges from mild confusion to deep coma, and it is often reversible when the underlying cause is treated. This chapter explains the common causes, features, and management of encephalopathy.

\section*{Epidemiology}
Encephalopathy is common in hospital settings, where it is often called delirium when it develops acutely. Hepatic encephalopathy affects a large proportion of people with advanced liver disease, while other forms complicate kidney failure, sepsis, and critical illness. Older people are especially vulnerable. Because the underlying conditions are common and growing, encephalopathy is a frequent reason for hospital admission and is associated with significant illness and mortality.

\section*{Aetiology and Pathophysiology}
The many causes of encephalopathy all share the effect of disturbing normal brain function.
\begin{itemize}
    \item \textbf{Hepatic encephalopathy:} the liver fails to clear toxins such as ammonia, which accumulate and affect the brain
    \item \textbf{Uraemic encephalopathy:} toxins accumulate in kidney failure
    \item \textbf{Hypoxic encephalopathy:} the brain is starved of oxygen after cardiac arrest, near-drowning, or severe breathing failure
    \item \textbf{Metabolic causes:} very low or high blood sugar, thyroid problems, and electrolyte disturbances
    \item \textbf{Infections:} sepsis and infections of the brain itself
    \item \textbf{Toxic causes:} alcohol, drugs, poisons, and medicines
    \item \textbf{Other causes:} Wernicke's encephalopathy from thiamine deficiency, and the brain changes of critical illness
    \item \textbf{Mechanism:} the underlying disturbance alters the chemical environment of the brain, impairing nerve cell function
\end{itemize}

\section*{Clinical Features}
The features of encephalopathy vary with the cause and severity.
\begin{itemize}
    \item Confusion, disorientation, and difficulty concentrating
    \item Drowsiness, lethargy, and reduced alertness
    \item Changes in behaviour and personality
    \item Slurred speech and unsteadiness
    \item Seizures in some forms
    \item Tremor, particularly a characteristic flapping tremor in hepatic encephalopathy
    \item Reduced consciousness, progressing to coma in severe cases
    \item In chronic forms: gradual cognitive decline
\end{itemize}

\section*{Differential Diagnosis}
Encephalopathy must be distinguished from other causes of confusion and altered consciousness.
\begin{itemize}
    \item Delirium from infection, pain, or medicines in hospital
    \item Dementia, which is chronic and progressive
    \item Stroke and brain injury
    \item Meningitis and encephalitis
    \item Psychiatric conditions, including psychosis and severe depression
    \item Seizures and non-convulsive epilepsy
\end{itemize}

\section*{When to Seek Medical Care}
Any sudden change in mental state, confusion, drowsiness, or altered consciousness requires urgent medical assessment. Families and carers should not assume changes are just "old age" or tiredness.

\textbf{Contact your local emergency services for:}
\begin{itemize}
    \item Sudden confusion or reduced consciousness
    \item Seizures or fits
    \item Difficulty breathing or blue lips
    \item Collapse or unresponsiveness
    \item Sudden weakness or facial drooping, which may indicate stroke
\end{itemize}

\section*{Diagnosis and Investigations}
Finding the cause of encephalopathy is the priority.
\begin{itemize}
    \item \textbf{History:} the onset and pattern of symptoms, medicines, and underlying conditions
    \item \textbf{Examination:} neurological and general examination
    \item \textbf{Blood tests:} liver and kidney function, blood sugar, electrolytes, oxygen, and infection markers
    \item \textbf{Imaging:} CT or MRI of the brain
    \item \textbf{EEG:} to assess brain activity and exclude seizures
    \item \textbf{Lumbar puncture:} when infection is suspected
    \item \textbf{Medicine review:} identifying drugs that impair brain function
\end{itemize}

\section*{Management}
Management targets the underlying cause while supporting the brain.
\begin{itemize}
    \item \textbf{Hepatic encephalopathy:} medicines that reduce ammonia, along with treatment of the trigger, such as infection or bleeding
    \item \textbf{Uraemic encephalopathy:} dialysis
    \item \textbf{Metabolic causes:} correction of glucose, electrolytes, and hormones
    \item \textbf{Oxygen and circulation:} supporting breathing and blood flow
    \item \textbf{Treating infection:} antibiotics for sepsis
    \item \textbf{Stopping toxins:} withdrawing alcohol and harmful medicines, and treating withdrawal
    \item \textbf{Thiamine:} given urgently when Wernicke's encephalopathy is suspected
    \item \textbf{Supportive care:} a calm environment, orientation, and nursing care
\end{itemize}

\section*{Complications}
\begin{itemize}
    \item Progression to coma
    \item Brain damage from prolonged oxygen deprivation
    \item Aspiration and pneumonia
    \item Falls and injury from confusion
    \item Pressure sores and blood clots from immobility
    \item Permanent cognitive impairment in some forms
    \item Death in severe or untreated cases
\end{itemize}

\section*{Prognosis and Outcomes}
The outlook depends on the cause and how quickly it is treated. Metabolic, toxic, and infectious causes often reverse fully when the underlying problem is corrected, and most people with hepatic encephalopathy recover with treatment. Hypoxic encephalopathy after cardiac arrest has a more guarded outlook. Early recognition and treatment give the best chance of full recovery.

\section*{Prevention and Self-Care}
\begin{itemize}
    \item Manage underlying conditions such as liver and kidney disease with regular care
    \item Avoid alcohol, which damages the liver and brain
    \item Take medicines as prescribed and review them with a doctor
    \item Seek urgent care for any sudden confusion or altered consciousness
    \item In hospital, help orient confused relatives with familiar items and company
    \item Follow rehabilitation advice after recovery
\end{itemize}

\section*{Sources and Further Reading}
The following reputable sources provide further information for healthcare students and patients seeking reliable, current advice about encephalopathy.
\begin{itemize}
    \item Healthdirect Australia. \textit{Encephalopathy and delirium information.} https://www.healthdirect.gov.au/
    \item Gastroenterology Society of Australia. \textit{Hepatic encephalopathy information.} https://www.gesa.org.au/
    \item NHS. \textit{Delirium and confusion in hospital.} https://www.nhs.uk/conditions/delirium/
    \item Mayo Clinic. \textit{Encephalopathy: overview.} https://www.mayoclinic.org/diseases-conditions/encephalitis/
    \item MedlinePlus. \textit{Encephalopathy.} https://medlineplus.gov/encephalopathy.html
    \item MSD Manual. \textit{Encephalopathy: professional overview.} https://www.msdmanuals.com/
\end{itemize}
